\documentclass[9pt,notes=hide]{beamer}

\usepackage{../../latex/custom}
\usepackage{../../latex/jsr}
\usepackage{../../latex/nsi}

\usepackage{animate}

\usetheme[progressbar=frametitle,block=fill]{metropolis}
\usepackage{appendixnumberbeamer}
\setbeamertemplate{frametitle continuation}{\hfill(\insertcontinuationcount)} %

\newcommand{\setsection}[1]{
\setbeamertemplate{footline}{
  \begin{beamercolorbox}[wd=\textwidth, sep=3ex]{footline}%
    #1
    \usebeamerfont{page number in head/foot}%
    \usebeamertemplate*{frame footer}
    \hfill%
    \usebeamertemplate*{frame numbering}
  \end{beamercolorbox}%
}
}

\usepackage{csquotes}
\usepackage{tikz}
\usetikzlibrary{automata, positioning, arrows, patterns}
\usetikzlibrary{matrix,decorations.pathreplacing,chains,fit,shapes,calc} %
\usetikzlibrary{calc,decorations.pathmorphing,decorations.markings}

\usepackage{pgfplots}
    \pgfplotsset{compat=1.11}
\usepackage{pgfplotsthemetol} %
\usepackage{subcaption} %
\usepackage{array}

\newcommand{\abs}[1]{\left| #1 \right|} %
\newcommand{\Init}{\mathrm{Init}} %
\newcommand{\comment}[1]{{\color{red}\textbf{{#1}}}}

\usepackage[super]{nth}

\newcommand{\expe}[2]{\la #1, #2 \ra}
\newcommand{\piB}{\pi_{\Image(B)^\perp}}
\newcommand{\piBTr}{\pi_{\Image(B)^\perp}^\Tr}
\DeclareMathOperator{\gauge}{g}
\DeclareMathOperator{\suppfun}{h}

\newcommand\Wider[2][3em]{%
\makebox[\linewidth][c]{%
  \begin{minipage}{\dimexpr\textwidth+#1\relax}
  \raggedright#2
  \end{minipage}%
  }%
}


\title{Piecewise semi-ellipsoidal control invariant sets}
\subtitle{59th Conference on Decision and Control}
\date{December 15, 2020}
\author{Beno\^{i}t \textsc{Legat}$^\dagger$, Sa\v{s}a V. \textsc{Rakovi\'c}$^*$, Rapha\"el M. \textsc{Jungers}$^\dagger$}
\institute{${}^\dagger$ UCLouvain, Belgium, ${}^*$ Beijing Institute of Technology, China}

\begin{document}

\begin{frame}%
  \maketitle
\end{frame}

\begin{frame}{Ellipsoidal invariant sets through gauge function}
  \begin{columns}
    \begin{column}{0.5\textwidth}
      \begin{block}{\alert{Gauge} function}
        $\mathcal{S} = \{\, x \mid \gauge(\mathcal{S}, x) \le 1 \,\}$.

        Inclusion: $\mathcal{S} \subseteq \mathcal{T} \Leftrightarrow \forall x, \gauge(\mathcal{S}, x) \ge \gauge(\mathcal{T}, x)$.

        Linear \alert{preimage}: $\gauge(A^{-1}\mathcal{S}, x) = \gauge(\mathcal{S}, Ax)$.

        \alert{Ellipsoid}: $\gauge(\mathcal{S}, x) = \sqrt{x^\Tr Q x}$.
      \end{block}
    \end{column}
    \begin{column}{0.5\textwidth}
      \begin{block}{Autonomous system}
        Classical LMI for stability:
        \begin{align*}
          x_{k+1} & = Ax_k\\
          A\mathcal{S} & \subseteq \mathcal{S}\\
          \mathcal{S} & \subseteq A^{-1}\mathcal{S}\\
          \gauge(\mathcal{S}, x) & \ge \gauge(A^{-1}\mathcal{S}, x)\\
          \gauge(\mathcal{S}, x) & \ge \gauge(\mathcal{S}, Ax)\\
          Q & \succeq A^\Tr Q A.
        \end{align*}
      \end{block}
    \end{column}
  \end{columns}
\end{frame}
\begin{frame}{Control invariant set: Algebraic approach}
  \begin{columns}
    \begin{column}{0.45\textwidth}
      \begin{block}{Discrete-time control system}
        \[ x_{k+1} = Ax_k + Bu_k. \]
      \end{block}

        Controlled invariance of $\mathcal{S}$:
        \[ \forall x \in \mathcal{S}, \exists u, A x + B u \in \mathcal{S}. \]

        \begin{block}{Linear state feedback}
          $u(x) = Kx$ $\to$ autonomous system:

          \vspace{-2em}
          \[
            x_{k+1} = (A + BK) x_k.
          \]
          \vspace{-2em}

          From previous slide:

          \vspace{-2em}
          \[
            \alert{Q} \succeq (A + B\alert{K})^\Tr \alert{Q} (A + B\alert{K}).
          \]
          \vspace{-2em}

          \alert{Cubic} matrix inequality!
        \end{block}
    \end{column}
    \begin{column}{0.55\textwidth}
      \begin{block}{Schur complement}

        \[
          \begin{bmatrix}
            \alert{Q} & (A + B\alert{K})^\Tr \alert{Q}\\
            \alert{Q}(A + B\alert{K}) & \alert{Q}
          \end{bmatrix} \succeq 0.
        \]

        \alert{Bilinear} matrix inequality!
      \end{block}
      \begin{block}{Change of variables}
        $Q^{-1}$ and $Y := KQ^{-1}$.

        \[
          \begin{bmatrix}
            \alert{Q^{-1}} & \alert{Q^{-1}}A^\Tr + \alert{Y}^\Tr B^\Tr\\
            A\alert{Q^{-1}} + B\alert{Y} & \alert{Q^{-1}}
          \end{bmatrix} \succeq 0.
        \]
        \alert{Linear} matrix inequality!
      \end{block}
    \end{column}
  \end{columns}
  \begin{center}
    \emph{How to generalize this for other sets ?}
  \end{center}
\end{frame}


\begin{frame}{Need for piecewise}
  \animategraphics[loop,controls,width=\linewidth]{5}{tow/0000}{000}{038}
\end{frame}

\begin{frame}{Piecewise semi-ellipsoids and polarity}
  \begin{columns}
    \begin{column}{0.48\textwidth}
      \begin{center}
        \begin{tikzpicture}[scale=2]
          \draw[thick]
          (-1, 0.02) to (-1, 0) to (0, -1)
          plot[domain=0:1, samples=40] ({\x}, {\x/2 - sqrt(4 - 3 * \x^2) / 2} )
          (1, 0) to (1, 1) to (0, 1)
          plot[domain=0:1, samples=40] ({-\x}, {sqrt(1 - \x^2)} );
          \draw[thick, densely dotted, ->] (0, -1) to (0.1, -1.2);
          \draw[thick, densely dotted, ->] (1, 0) to (1.2, -0.1);
        \end{tikzpicture}
      \end{center}
      \begin{equation*}
        \begin{cases}
          x^2 & \text{ if } 0 \le y \le x,\\
          y^2 & \text{ if } 0 \le x \le y,\\
          x^2 + y^2 & \text{ if } x \le 0 \le y,\\
          (x + y)^2 & \text{ if } x, y \le 0,\\
          x^2 - xy + y^2 & \text{ if } y \le 0 \le x.
        \end{cases}
      \end{equation*}
    \end{column}
    \begin{column}{0.52\textwidth}
      \begin{center}
        \begin{tikzpicture}[scale=2]
          \draw[thick] %
          (-1, 0.02) to (-1, -1) to (0.5, -1)
          plot[thick, domain=0.5:1, samples=40] ({\x}, {-\x/2 - sqrt(3 - 3 * \x^2) / 2} )
          (1, -0.52) to (1, 0) to (0, 1)
          plot[thick, domain=0:1, samples=40] ({-\x}, {sqrt(1 - \x^2)} );
          \draw[thick, densely dotted, ->] (0, 0) to (0.5, -1);
          \draw[thick, densely dotted, ->] (0, 0) to (1, -0.5);
        \end{tikzpicture}
      \end{center}
      \begin{equation*}
        \begin{cases}
          (x + y)^2 & \text{ if } 0 \le x, y,\\
          x^2 + y^2 & \text{ if } x \le 0 \le y,\\
          x^2 & \text{ if } x \le y \le 0,\\
          y^2 & \text{ if } y \le x, 2x \le -y,\\
          \dfrac{4}{3}(x^2 + xy + y^2) & \text{ if } 2x \ge -y, x \le -2y,\\
          x^2 & \text{ if } x \ge -2y, y \ge 0.
        \end{cases}
      \end{equation*}
    \end{column}
  \end{columns}
\end{frame}

\begin{frame}{Control invariance and algebraic invariance through support function}
  \begin{columns}
    \begin{column}{0.5\textwidth}
      \begin{block}{\alert{Support} function}
        \begin{align*}
          \suppfun(\mathcal{S}, y) & = \sup_{x \in \mathcal{S}} y^\Tr x\\
          & = \gauge(\polar{\mathcal{S}}, x)
        \end{align*}

        Inclusion: $\mathcal{S} \subseteq \mathcal{T} \Leftrightarrow \forall y, \suppfun(\mathcal{S}, y) \le \suppfun(\mathcal{T}, y)$.

        Linear \alert{image}: $\suppfun(A\mathcal{S}, y) = \suppfun(\mathcal{S}, A^\Tr y)$.

        \alert{Ellipsoid}:
        \begin{align*}
          \gauge(\mathcal{S}, x) & = \sqrt{x^\Tr Q x}\\
          \suppfun(\mathcal{S}, y) & = \sqrt{y^\Tr Q^{-1} y}
        \end{align*}
      \end{block}
    \end{column}
    \begin{column}{0.5\textwidth}
      \begin{block}{\alert{Control} system}
        \[ x_{k+1} = Ax_k + Bu_k. \]
      \end{block}

      Controlled invariance of $\mathcal{S}$:
      \[ \forall x \in \mathcal{S}, \exists u, A x + B u \in \mathcal{S}. \]

      \vspace{-1em}

      Equivalent to
      \begin{align*}
        \piB{} \mathcal{S} & \subseteq \piB{} A \mathcal{S}\\
        \suppfun(\mathcal{S}, \piBTr{} y) & \le \suppfun(\mathcal{S}, A^\Tr \piBTr{} y)
      \end{align*}

      \begin{block}{\alert{Algebraic} system}
        \[ \piB{} x_{k+1} = \piB{} Ax_k. \]
      \end{block}
    \end{column}
  \end{columns}
  \[
    \piB{}Q^{-1}\piBTr{} \preceq  \piB{} A Q^{-1} A^\Tr \piBTr{}
  \]
\end{frame}

\begin{frame}{LMI formulation}

  \begin{block}{Piecewise semi-ellipsoid}
    Fixed polyhedral conic partition $\mathcal{P}_i$:
    \[
      \suppfun(\mathcal{S}, x) =
        \sqrt{x^\Tr Q_i x} \quad \text{ if } x \in \mathcal{P}_i
    \]
  \end{block}

  \begin{align*}
    \text{find}
    & \quad Q_i \succeq 0, u_{i, j} \in \R^n \quad \text{ such that}\\
      x^\Tr C Q_i C^\Tr x & \leq x^\Tr E Q_j E^\Tr x,
      \forall x \in C^{-\Tr} \mathcal{P}_i \cap E^{-\Tr} \mathcal{P}_j,\\
      & \quad \qquad \qquad \qquad \forall i, j \in \{1, \ldots, m\}\\
      v^\Tr Q_i v & \le 1, \quad \forall v \in \Vrep(\polar{\mathcal{X}}) \cap \mathcal{P}_i, \forall i \\
      Q_i - Q_j & = u_{i,j} n_{ij}^\Tr + n_{ij} u_{i,j}^\Tr,
      \quad \forall i, j \in \mathcal{N}\\
      Q_j n_{ij} - Q_i n_{ij} & \in \polar{(\mathcal{P}_i \cap \mathcal{P}_j)}, \quad \forall i, j \in \mathcal{N}.
  \end{align*}
\end{frame}

\begin{frame}{Example}
  \begin{columns}
    \begin{column}{0.5\textwidth}
      \alert{Primal} space
      \begin{center}
        \includegraphics[trim=3.7cm 0cm 3.5cm 0cm, clip, width=0.65\textwidth]{SimpleControlledInvariantEllipsoid.png}
        \includegraphics[trim=3.7cm 0cm 3.5cm 0cm, clip, width=0.65\textwidth]{SimpleControlledInvariantPiecewiseSemiEllipsoid.png}
      \end{center}
    \end{column}
    \begin{column}{0.5\textwidth}
      \alert{Dual} space
      \begin{center}
        \includegraphics[trim=3.7cm 0cm 3.5cm 0cm, clip, width=0.65\textwidth]{SimpleControlledInvariantPolarEllipsoid.png}
        \includegraphics[trim=3.7cm 0cm 3.5cm 0cm, clip, width=0.65\textwidth]{SimpleControlledInvariantPolarPiecewiseSemiEllipsoid.png}
      \end{center}
    \end{column}
  \end{columns}
\end{frame}

\begin{frame}{Conclusion}
  \begin{itemize}
    \item Ellipsoids are \alert{cheap} to compute but conservative.
    \item Polyhedra representation are \alert{expressive} but can have huge size in some regions.
    \item Piecewise semi-ellipsoids combine \alert{advantages} of both.
  \end{itemize}
  Future work:
  \begin{itemize}
    \item How to obtain conic partition ?
      \begin{itemize}
        \item A few backward propagation of a polytope.
        \item Refine with dual solution.
      \end{itemize}
    \item Piecewise polysets (polynomial support function).
  \end{itemize}
\end{frame}







\end{document}
